%第1章





\begin{example}{牛顿第二定律与微积分结合*}
    质点从原点出发，沿$x$轴正向运动，初速度为$\SI{2}{m/s}$，加速度与位置的关系是$a=4+4x$(SI). (1) 求质点速率$v$与$x$的关系; (2) 求$x$与$t$的关系。
    \tcblower
    {\fangsong 思路：若已知$x$与$a$的关系，需利用$a=\dfrac{dv}{dt}=\dfrac{dv}{dx}\dfrac{dx}{dt}=v\dfrac{dv}{dx}$进行代换处理。}
    \tcbline
    (1)~利用$a=\dfrac{dv}{dt}=\dfrac{dv}{dx}\dfrac{dx}{dt}=v\dfrac{dv}{dx}$，结合题意可得$v\dfrac{dv}{dx}=4+4x$

    \vspace{0.5em}
    
    分离变量得到$vdv=(4+4x)dx$，两边同时积分$\int{vdv}=\int{(4+4x)dx}$，得$\dfrac{1}{2}v^2=4x+2x^2+C$

    \vspace{0.25em}

    初始条件为$v=2,x=0$，代入求得$C=2$.
    
    结合题意，质点正向运动，所以$v>0$，因此整理得$v=\sqrt{2(2x^2+4x+2)}=2(x+1)$

    \vspace{0.25em}

    (2)~根据上一问结果，$v=\dfrac{dx}{dt}=2(x+1)$，分离变量得$\dfrac{dx}{x+1}=2dt$

    \vspace{0.5em}

    两边同时积分$\int{\dfrac{dx}{x+1}}=\int{2dt}$，得$\ln{(x+1)}=2t+C$

    \vspace{0.25em}

    初始条件为$t=0$时，$x=0$，代入求得$C=0$.~整理得$x=e^{2t}-1$
    
\end{example}

\begin{example}{一维运动的物理量}
    质点沿$x$轴运动，位置与时间的关系是$x=3t-t^3+6$(SI)。
    
    (1)~求任意时刻的速度和加速度；(2)~判断什么时候向正方向运动，什么时候向负方向运动；(3)~判断什么时候加速、什么时候减速；(4)~求前$\SI{2}{s}$内的路程、位移、平均速度、平均速率。
    \tcblower
    {\fangsong 思路：对$x$求两次导，分别可得速度、加速度与时间的关系，再结合实际意义进行分析。}
    \tcbline

    \setlength{\parindent}{2em}
    
    \noindent (1)~~将$x$对$t$求导，可得$v=x^{\prime}=3-3t^2$，再次求导可得$a=x^{\prime\prime}=v^{\prime}=-6t$
    \note{在一维问题中，矢量运算退化成标量运算，可省略上方的箭头，用正负表示其方向。}

    \noindent (2)~~当$v>0$时，即$0\leq t<\SI{1}{s}$时，向正方向运动；而$v<0$时，即$t>\SI{1}{s}$时，向负方向运动。

    \noindent (3)~~因为$a=-6t<0$，因此当$0\leq t<\SI{1}{s}$时，$v>0$，加速度与速度反向，为减速；
    
    而$t>\SI{1}{s}$时，$v<0$，加速度与速度同向，为加速。

    \noindent (4)~~当$t=0,\SI{1}{s},\SI{2}{s}$时，代入表达式可求出$x=\SI{6}{m},\SI{8}{m},\SI{4}{m}$

    位移就是位置变化量，只看初末状态即可，$\Delta x=\SI{4}{m}-\SI{6}{m}=\SI{-2}{m}$

    路程需要按转向时刻分段求解，第$\SI{1}{s}$路程为$\SI{2}{m}$，第$\SI{2}{s}$路程为$\SI{4}{m}$，因此总路程为$\SI{6}{m}$

    平均速度等于位移除以时间，等于$\SI{-1}{m/s}$；平均速率等于路程除以时间，等于$\SI{3}{m/s}$.

    \tcbline
    
\end{example}





%第2章
\begin{minipage}{0.78\linewidth}
    \setlength{\parskip}{0.5em}
    \linespread{1.3}
    【补充练习*】半径为$R$的圆柱形桶中装有一些水，桶绕$z$轴以恒定角速度$\omega$转动，求稳定时的水面形状。（提示：取水面上的流体微团进行受力分析，注意该处水面倾角$\theta$与$z,r$坐标的关系。）
    \end{minipage}
    \hfill
    \begin{minipage}{0.18\textwidth}
        \centering \includegraphics[width=\textwidth]{asset/ch2/T1.png}
    \end{minipage}

%第4章表格
\subsection{质点与刚体的对比}

\vspace{1em}

\begin{table}[htb]
    \centering
    \begin{tabular}{|c|c|c|}
        \hline
         & 质点 & 刚体 \\
        \hline
        运动的描述 & 线量$x,v,a$ & 角量$\theta, \omega, \alpha$ \\
        \hline
        惯性的度量 & 质量$m$ & 转动惯量$J$ \\
        \hline
        运动定律 & $F=ma$ & $M=J\alpha$ \\
        \hline
        \rule{0pt}{18pt}\rule[-12pt]{0pt}{0pt} 动能 & $E_k=\dfrac{1}{2}mv^2$ & $E_k=\dfrac{1}{2}J\omega^2$ \\
        \hline
        \rule{0pt}{18pt}\rule[-12pt]{0pt}{0pt} 动能定理 & $\displaystyle\int{Fdx}=\Delta E_k$ & $\displaystyle\int{Md\theta}=\Delta E_k$ \\
        \hline
        重力势能 & $E_p=mgh$ & $E_p=mgh_C$ \\
        \hline
        机械能守恒条件 & 只有保守内力做功 & 只有保守内力做功 \\
        \hline
        动量/角动量 & $p=mv$ & $L=J\omega$ \\
        \hline
        \rule{0pt}{18pt}\rule[-12pt]{0pt}{0pt} 动量/角动量定理 & $\displaystyle\int{Fdt}=\Delta p$ & $\displaystyle\int{Mdt}=\Delta L$ \\
        \hline
        动量/角动量守恒条件 & 合外力为零 & 合外力矩为零 \\
        \hline
    \end{tabular}
\end{table}